\\newpage
\\section{模型评价}

\\subsection{预测模型评价}
\\begin{itemize}
    \\item \\textbf{精度优势}：Transformer模型MAPE=3.8\\%，优于ARIMA（5.3\\%）约28\\%
    \\item \\textbf{泛化能力}：在测试集上表现稳定，MAPE波动不超过0.5个百分点
    \\item \\textbf{计算效率}：单步预测耗时约0.02秒（CPU），满足实时调度需求
\\end{itemize}

\\subsection{调度模型评价}
\\begin{itemize}
    \\item \\textbf{收敛性}：Q-Learning在300轮后reward曲线趋于平稳
    \\item \\textbf{Pareto质量}：NSGA-II求解的解集覆盖完整的Pareto前沿
    \\item \\textbf{对比优势}：相较于基准规则调度，综合指标提升约12\\%
\\end{itemize}

\\subsection{模型优缺点}
\\textbf{优点}：
\\begin{itemize}
    \\item 预测模型融合STL分解与Attention机制，精度高
    \\item 调度模型考虑多目标冲突，提供多种可行方案
    \\item 框架模块化，易于扩展
\\end{itemize}
\\textbf{缺点}：
\\begin{itemize}
    \\item 纯NumPy实现，计算效率有待提升
    \\item 状态空间离散化导致精度损失
    \\item 未考虑极端天气等扰动因素
\\end{itemize}
